\chapter
[\CO{[WIS]} Wissenschaft und Forschung]
{\CC{[WIS]} Wissenschaft und Forschung}
\label{chp:WIS}

\minitoc

\section{Von der Studie zur Lehrmeinung}

\subsection{Grundlagen -- Was ist Wissenschaft?}

\subsubsection{Wissensgewinn: Am Anfang war die Frage -- am
Ende auch}

\subsubsection{Welt in Zahlen ausgedrückt: Die Rolle der
Statistik}

\subsubsection{Was sind Studien?}

\subsubsection{Fachgesellschaften}

In der Welt der Wissenschaft enstehen und entstanden sog. Fachgesellschaften,
die sich bestimmten Fachgebieten, Erkrankungen und Fragestellungen widmen.

\cite{Bahr200701}: \bibentry{Bahr200701}

\Layout{0}{Medizinische Fachgesellschaften}
{%
\label{topic:fachgesellschaft}
\index{Fachgesellschaft}
Die Tätigkeit von medizinische Fachgesellschaften umfasst oft die Forschung,
Fortbildung und die Herausgabe von
\EE{Empfehlungen} und Leitlinien zur Diagnostik, Therapie und Vorbeugung
u.\,ä. Diese Empfehlungen sind nicht "`gesetzlich bindend"', aber geben
i.\,d.\,R. den \EE{Stand der Wissenschaft} wieder. Ausserdem müssen diese, oft
allgemein formulierten, Empfehlungen an die örtlichen Gegebenheiten angepasst
werden. Verschiedene Fachgesellschaften können -- auch zu den gleichen Themen --
unterscheidliche Meinungen haben und unterschiedliche Empfehlungen aussprechen.
}{%
\begin{itemize}
   \item Forschung, Fortbildung 
   \item Herausgabe von Empfehlungen für Diagnostik, Therapie und
   Vorbeugung (u.\,a.)
   \item Nicht bindend, aber "`Stand der Wissenschaft"'
\end{itemize}
}{}{}{}


\Layout{0}{ERC}
{%
\label{topic:erc}
\index{ERC}
\index{Fachgesellschaft!ERC}
Eine wichtige Fachgesellschaft ist das \T{ERC} (\E{European Resuscitation
Council}). Bekannt ist das \Abk{ERC} u.a. durch seine Leitlinien zur
Wiederbelebung. Alle 5 Jahre werden diese
"`Guidelines\footnote{\Lang{engl.} \emph{Guideline}: Richtlinie,
Empfehlung}"' aktualisiert. Jeder neuen Auflage der Richtlinien geht ein aufwendiger Forschungs- und Diskussionsprozess voraus: Es werden hunderte Studien mit vielen verschiedenen Fragestellungen an vielen
Forschungszentren auf der ganzen Welt durchgeführt, die Ergebnisse ausgewertet und anschließend diskutiert. 

Die derzeit aktuelle Version
der ERC-Guidelines kam im Jahr 2005 heraus("`\E{ERC 2005}"'), die nächste
Aktualisierung ist mit Ende des Jahre 2010 zu erwarten.
}{%
\begin{itemize}
   \item Leitlinien zur Reanimation
   \item Aktualisierung alle 5 Jahre
   \item Letzte 2005, nächste Ende 2010
\end{itemize}

}{}{}{}


\section{Statistik}


\Captionof{table}{Statistische Tests mit lustigen Namen -- Eine Übersicht.}